% ----- CHAUPAI 35 -----

\newpage

\subsection*{\deva{पञ्चत्रिंशत्तम चौपाई} (Thirty-Fifth Chaupai)}
\addcontentsline{toc}{subsection}{\deva{पञ्चत्रिंशत्तम चौपाई} (Thirty-Fifth Chaupai) :   \deva{और देवता...}}

\begin{minipage}[t]{0.55\textwidth}
\vspace{0pt}

{\large \linespread{1.5}\selectfont
\deva{और देवता चित्त न धरई।}\\
\deva{हनुमत सेइ सर्ब सुख करई॥}
\par}

\vspace{0.5em}

{\large \linespread{1.5}\selectfont
\telu{ఔర దేవతా చిత్త న ధరఈ।}\\
\telu{హనుమత సేఇ సర్బ సుఖ కరఈ॥}
\par}

\vspace{0.5em}

{\large \linespread{1.3}\selectfont
\textit{aura devatā citta na dharaī |}\\
\textit{hanumata sei sarba sukha karaī ||}
\par}
\end{minipage}%
\hfill
\begin{minipage}[t]{0.42\textwidth}
\vspace{0pt}
\begin{tcolorbox}[anchorbox]
\textbf{One-Pointed Focus:} A magnifying glass only starts a fire if held completely still. True academic or career success demands absolute, laser-focused dedication to one primary goal (\textit{Ananya Bhakti}) rather than constantly bouncing between different distractions.
\vspace{0.5em}\newline He ruthlessly blocked out all other distractions, organizing massive efforts with absolute, single-minded focus on his goal.\newline\null\hfill\href{https://www.valmiki.iitk.ac.in/sloka?field_kanda_tid=4&language=dv&field_sarga_value=39}{\textit{Kishkindha Kanda, 39:10}}
\end{tcolorbox}
\end{minipage}

\vspace{0.5em}
\subsubsection*{\textbf{Visual Rhythm Map:}}
\begin{table}[h!]
\centering
\renewcommand{\arraystretch}{1.2}
\setlength{\tabcolsep}{0pt}
\begin{tabularx}{\textwidth}{|*{16}{>{\centering\arraybackslash\hsize=1.0\hsize}X|}}
\hline
\multicolumn{3}{|>{\centering\arraybackslash\hsize=3.0\hsize}X|}{\textbf{\guru{औ}\laghu{र}} \par (3)} &
\multicolumn{5}{>{\centering\arraybackslash\hsize=5.0\hsize}X|}{\textbf{\guru{दे}\laghu{व}\guru{ता}} \par (5)} &
\multicolumn{3}{>{\centering\arraybackslash\hsize=3.0\hsize}X|}{\textbf{\guru{चि}\laghu{त्त}} \par (3)} &
\multicolumn{1}{>{\centering\arraybackslash\hsize=1.0\hsize}X|}{\textbf{\laghu{न}} \par (1)} &
\multicolumn{4}{>{\centering\arraybackslash\hsize=4.0\hsize}X|}{\textbf{\laghu{ध}\laghu{र}\guru{ई}} \par (4)} \\
\hline
\end{tabularx}
\vspace{-1.5em}
\end{table}

\begin{table}[h!]
\centering
\renewcommand{\arraystretch}{1.2}
\setlength{\tabcolsep}{0pt}
\begin{tabularx}{\textwidth}{|*{16}{>{\centering\arraybackslash\hsize=1.0\hsize}X|}}
\hline
\multicolumn{4}{|>{\centering\arraybackslash\hsize=4.0\hsize}X|}{\textbf{\laghu{ह}\laghu{नु}\laghu{म}\laghu{त}} \par (4)} &
\multicolumn{3}{>{\centering\arraybackslash\hsize=3.0\hsize}X|}{\textbf{\guru{से}\laghu{इ}} \par (3)} &
\multicolumn{3}{>{\centering\arraybackslash\hsize=3.0\hsize}X|}{\textbf{\guru{स}\laghu{र्ब}} \par (3)} &
\multicolumn{2}{>{\centering\arraybackslash\hsize=2.0\hsize}X|}{\textbf{\laghu{सु}\laghu{ख}} \par (2)} &
\multicolumn{4}{>{\centering\arraybackslash\hsize=4.0\hsize}X|}{\textbf{\laghu{क}\laghu{र}\guru{ई}} \par (4)} \\
\hline
\end{tabularx}
\vspace{-1.5em}
\end{table}

\vspace{1em}
\textbf{ \deva{सरल-संस्कृतम्}:}\\
 \deva{मनसि कस्यापि अन्यस्य देवतायाः ध्यानं न धर्तव्यम्।}\\
 \deva{केवलं हनुमतः सेवया (भक्त्या) एव सर्वाणि सुखानि प्राप्यन्ते॥}\\


Let not the mind hold or rely on any other deity.\\
By serving Hanuman alone, one attains all possible happiness.


\vspace{1em}
\newpage
\textbf{\deva{प्रतिपदार्थः} (Meaning of each word):}
\begin{table}[h!]
\centering
\renewcommand{\arraystretch}{1.2}
\begin{tabularx}{\textwidth}{@{} l l X X @{}}
\toprule
\textbf{Awadhi} & \textbf{Sanskrit} & \textbf{English} & \textbf{Grammar \& Notes} \\
\midrule
\deva{और} / \telu{ఔర}  &  \deva{अन्य / परः}  &  Other  &   \\
\deva{देवता} / \telu{దేవతా}  &  \deva{देवता}  &  Deity / Gods  &   \\
\deva{चित्त} / \telu{చిత్త}  &  \deva{चित्ते / मनसि}  &  In the mind  &   \\
\deva{न धरई} / \telu{న ధరఈ}  &  \deva{न धरति}  &  Does not hold/keep  & `y` $\rightarrow$ `ee` verb \\ 
\deva{हनुमंत} / \telu{హనుమంత}  &  \deva{हनुमन्तम्}  &  Hanuman  &   \\
\deva{सेइ} / \telu{సేఇ}  &  \deva{सेवया / आराध्य}  &  By serving  & Absolutive `-i` ending \\
\deva{सर्ब सुख} \gotchaicon / \telu{సర్బ సుఖ}  &  \deva{सर्व-सुखानि}  &  All happiness  &  \\ 
\deva{करई} / \telu{కరఈ}  &  \deva{करोति}  &  Achieves / Does  & `y` $\rightarrow$ `ee` verb \\ 
\bottomrule
\end{tabularx}
\end{table}

\begin{tcolorbox}[poetrybox]
\textbf{\deva{सर्व} $\rightarrow$ \deva{सर्ब} \gotchaicon\ — The Invisible Case Agreement:}\\
Two shifts happen simultaneously in \deva{सर्ब सुख}. First, the familiar \textbf{\deva{व} $\rightarrow$ \deva{ब} shift} turns \deva{सर्व} into \deva{सर्ब}. 
Second, and this is the real gotcha, in classical Sanskrit, \deva{सर्व} must agree in gender, number, and case with the noun it modifies, if they appear as separate words instead of in a compund word like \deva{सर्वसुखसंपन्न}. 
Since \deva{सुख} (happiness) is a \textit{neuter plural}, Sanskrit demands \deva{सर्वाणि सुखानि}.
Awadhi simply ignores this entirely! \deva{सर्ब} is used as a frozen, indeclinable adjective, it never changes its form regardless of the noun's gender or case. 
This ''case-dropping'' is one of the most significant structural differences between Sanskrit and Awadhi.
\end{tcolorbox}

\newpage
